\documentclass[12pt,a4paper]{article}
\usepackage[margin=1in]{geometry}
\usepackage{amsmath,amssymb}
\usepackage{booktabs}
\usepackage{enumitem}
\usepackage{hyperref}

\title{Thiết kế lại mã hóa hai cấp cho GA + Tabu\\
(Truck Separators + Drone Assignment/Boundary)}
\author{Drone\_Resupply Project}
\date{\today}

\begin{document}
\maketitle

\section{Mục tiêu}
Tài liệu này đặc tả phương án mã hóa mới cho pipeline GA + Tabu.
Mục tiêu là chuyển từ chromosome chỉ gồm sequence sang biểu diễn hai cấp sao cho:
\begin{itemize}[leftmargin=1.2cm]
    \item tách khách hàng theo truck một cách tường minh ở cấp trên,
    \item điều khiển phục vụ bằng drone và biên chuyến (multi-trip) tường minh ở cấp dưới,
    \item tương thích với kiểm tra khả thi (capacity, endurance, time window) ở decoder.
\end{itemize}

\section{Mã hóa cuối cùng}
\subsection{Cấp trên: \texttt{sequence\_with\_truck\_sep}}
Chromosome cấp trên là một vector duy nhất gồm:
\begin{itemize}[leftmargin=1.2cm]
    \item ID khách hàng: $1..N$
    \item ID separator truck: $N+1..N+K-1$ (đúng $K-1$ separator)
\end{itemize}

Ý nghĩa:
\begin{itemize}[leftmargin=1.2cm]
    \item thứ tự khách = thứ tự phục vụ,
    \item separator chia khách thành các segment truck.
\end{itemize}

Ví dụ ($N=8, K=3$):
\[
\text{Upper}=[3,7,1,9,8,2,10,6,5,4]
\]
với separator là $9,10$.
Giải mã segment truck:
\begin{itemize}[leftmargin=1.2cm]
    \item Truck 0: $[3,7,1]$
    \item Truck 1: $[8,2]$
    \item Truck 2: $[6,5,4]$
\end{itemize}

\textbf{Quy tắc hợp lệ của upper}
\begin{itemize}[leftmargin=1.2cm]
    \item mỗi khách hàng xuất hiện đúng một lần,
    \item mỗi separator truck xuất hiện đúng một lần,
    \item không có ID thừa.
\end{itemize}

\subsection{Cấp dưới: \texttt{drone\_assign} + \texttt{boundary}}
Hai vector cấp dưới được đánh chỉ số theo \textbf{vị trí khách hàng} sau khi bỏ separator khỏi upper.

\begin{itemize}[leftmargin=1.2cm]
    \item \texttt{drone\_assign[i]}:
    \begin{itemize}
        \item $0$: khách $i$ do truck xử lý,
        \item $1..D$: khách $i$ được gán cho drone $d$.
    \end{itemize}
    \item \texttt{boundary[i]}:
    \begin{itemize}
        \item $1$: khách $i$ là điểm bắt đầu một chuyến drone mới,
        \item $0$: khách $i$ nối tiếp chuyến trước (cùng truck segment và cùng drone).
    \end{itemize}
\end{itemize}

\textbf{Quy tắc hợp lệ của lower}
\begin{itemize}[leftmargin=1.2cm]
    \item nếu \texttt{drone\_assign[i]}$=0$ thì ép \texttt{boundary[i]}$=0$,
    \item nếu \texttt{boundary[i]}$=0$ thì phải có khách drone hợp lệ trước đó để nối,
    \item không được nối chuyến qua separator truck,
    \item nếu đổi drone ID thì \texttt{boundary} bắt buộc bằng 1.
\end{itemize}

\section{Vì sao mô hình này bao quát hơn trong thực tế}
So với cách đặt separator drone ở upper, mô hình này ổn định hơn vì:
\begin{itemize}[leftmargin=1.2cm]
    \item phân tách truck được biểu diễn tường minh ở upper,
    \item phục vụ bằng drone được biểu diễn tường minh ở lower,
    \item nhiều chuyến của cùng một drone được biểu diễn bằng \texttt{boundary},
    \item không cần ngữ nghĩa separator drone cố định hoặc không giới hạn.
\end{itemize}

Lưu ý: tính khả thi vẫn do decoder kiểm tra (capacity/endurance/time-window), không hard-code hết trong chromosome.

\section{Chiến lược khởi tạo}
\subsection{Bước 1: Giữ nguyên bộ sinh upper hiện có}
Giữ nguyên 4 nhóm khởi tạo upper:
\begin{itemize}[leftmargin=1.2cm]
    \item 10\% Random
    \item 30\% Greedy-Time
    \item 30\% Sweep-Angle
    \item 30\% Nearest-Neighbor
\end{itemize}

\subsection{Bước 2: Chèn separator truck bằng chia đều}
Với sequence khách độ dài $N$ và $K$ truck:
\[
\text{base}=\left\lfloor\frac{N}{K}\right\rfloor,\quad
\text{rem}=N\bmod K
\]
Kích thước từng segment:
\[
\text{seg\_len}[i]=
\begin{cases}
\text{base}+1 & i<\text{rem}\\
\text{base} & \text{ngược lại}
\end{cases}
\]
Chèn separator sau mỗi segment, trừ segment cuối.

Cải tiến tùy chọn: cho separator trượt trong cửa sổ $\pm 2$ để chọn điểm cắt có chi phí cạnh cục bộ nhỏ hơn.

\subsection{Bước 3: Sinh lower bằng greedy decode nhẹ}
Chạy hàm nhẹ \texttt{greedyDecodeForInit(upperWithSep, data)} để tạo:
\begin{itemize}[leftmargin=1.2cm]
    \item \texttt{drone\_assign} trong miền $[0..D]$,
    \item \texttt{boundary} trong miền $\{0,1\}$.
\end{itemize}

Hàm init-light này nên:
\begin{itemize}[leftmargin=1.2cm]
    \item tách segment truck từ upper,
    \item gán truck vs drone theo rule greedy hiện có,
    \item tạo trip boundary,
    \item bỏ qua phần mô phỏng timeline nặng.
\end{itemize}

Chuẩn hóa sau cùng:
\begin{itemize}[leftmargin=1.2cm]
    \item nếu \texttt{drone\_assign[i]}$=0$ thì \texttt{boundary[i]}$\leftarrow 0$,
    \item boundary không được nối qua separator truck,
    \item khách không đủ điều kiện drone bị ép \texttt{drone\_assign=0}.
\end{itemize}

\section{Contract decoder (cho đánh giá GA/Tabu)}
Đầu vào decoder: $(\texttt{upper},\texttt{drone\_assign},\texttt{boundary})$.

Quy trình:
\begin{enumerate}[leftmargin=1.2cm]
    \item tách upper thành các segment truck,
    \item bỏ separator và dựng map vị trí khách,
    \item dựng các chuyến drone từ \texttt{drone\_assign}+\texttt{boundary} trong từng segment,
    \item mô phỏng hoạt động truck/drone và tính objective,
    \item áp kiểm tra khả thi và penalty.
\end{enumerate}

\section{Ghi chú triển khai cho GA + Tabu}
\subsection{Toán tử GA}
\begin{itemize}[leftmargin=1.2cm]
    \item Crossover upper: OX/PMX/CX trên upper đầy đủ + repair separator.
    \item Crossover lower: one-point hoặc uniform cho \texttt{drone\_assign}/\texttt{boundary}.
    \item Mutation upper: swap/inversion/displacement có repair nhận biết separator.
    \item Mutation lower: flip drone value, flip boundary, split/merge boundary block.
\end{itemize}

\subsection{Bộ nhớ Tabu và neighborhood}
\begin{itemize}[leftmargin=1.2cm]
    \item Tabu key phải hash cả 3 vector.
    \item Neighborhood nên chứa cả move upper và move lower.
    \item Cách thực dụng: xen kẽ vòng upper-step và lower-step.
\end{itemize}

\section{Ví dụ kiểm tra nhanh}
Giả sử upper sau khi bỏ separator có thứ tự khách:
\[
[3,7,1,8,2,6,5,4]
\]
Lower vectors:
\[
\texttt{drone\_assign}=[1,1,0,2,2,0,1,1]
\]
\[
\texttt{boundary}=[1,0,0,1,0,0,1,0]
\]
Diễn giải:
\begin{itemize}[leftmargin=1.2cm]
    \item Chuyến A (drone 1): khách 3,7
    \item khách 1 do truck phục vụ
    \item Chuyến B (drone 2): khách 8,2
    \item khách 6 do truck phục vụ
    \item Chuyến C (drone 1): khách 5,4
\end{itemize}

\section{Kết luận}
Phương án chốt:
\begin{itemize}[leftmargin=1.2cm]
    \item Upper: thứ tự khách + separator truck,
    \item Lower: drone assignment + trip boundary,
    \item Init: giữ 10/30/30/30 cho upper rồi chèn separator theo chia đều,
    \item Lower init: suy ra từ greedy decode nhẹ cho giai đoạn khởi tạo.
\end{itemize}
Thiết kế này giữ được phong cách code hiện tại và cho phép tối ưu đầy đủ theo 3 vector trong GA + Tabu.

\end{document}
